\documentclass{article}

%%% Package List

\usepackage{datetime} % Dates.
\usepackage{fancyhdr} % Headers.
\usepackage{lastpage} % Page References.

\usepackage{graphicx} % Inserting Images.
\usepackage{float} % Postitioning Images.

\usepackage{dirtytalk} % Speech & Quote Marks.
\usepackage{hyperref} % Hyperlinks and References.
\usepackage{xcolor} % Coloured Text.
\usepackage{enumitem} % Letter Bullets. 

\usepackage{amssymb} % Maths Symbols.
\usepackage{amsmath} % Maths Tools.

\usepackage{algpseudocode} % Pseudocode.
\usepackage{algorithm} % Algorithms.
\usepackage{minted} % Code Snippets.

\usepackage[margin=1in]{geometry} % Reduces the Page Margins.

%%%



%%% Custom Formats and Commands

\newcommand{\templateversion}{1.5.1} % Do not edit this variable.

\hypersetup{colorlinks=true, urlcolor=black, linkcolor=black}


% Date Formats.
\newdateformat{monthyeardate}
{%
    \monthname[\THEMONTH] \THEYEAR
}

\newdateformat{yeardate}
{%
    \THEYEAR
}


% Colouring and Spacing.
\newcommand{\colour}[2] {\color{#1}#2\color{black}}
\newcommand{\separate}[1] {\vspace{#1 pt}\noindent}


% Maths Number Sets.
\newcommand{\R}{\mathbb{R}}
\newcommand{\N}{\mathbb{N}}
\newcommand{\C}{\mathbb{C}}
\newcommand{\Q}{\mathbb{Q}}
\newcommand{\Z}{\mathbb{Z}}
\newcommand{\I}{\mathbb{I}}


% Custom Formats.
\newcommand{\definition}[2]{\textbf{Definition #1}: #2\label{def:#1}}
\newcommand{\example}[2]{\textbf{Example #1}: #2\label{ex:#1}}
\newcommand{\conjecture}[2]{\textbf{Conjecture #1}: #2\label{cnj:#1}}
\newcommand{\numberedentry}[3]{\textbf{#2 #1}: #3\label{entry:#1}}
\newcommand{\proof}[2]{\textit{Proof}: #2 \begin{flushright}\(\square\)\end{flushright}\label{prf:#1}}
\newcommand{\theorem}[3]{\textbf{Theorem #1}: #2 \par\separate{3}\proof{thm:#1}{#3}\label{thm:#1}}
\newcommand{\lemma}[3]{\textbf{Lemma #1}: #2 \par\separate{3}\proof{lem:#1}{#3}\label{lem:#1}}


% Join Symbols.
\newcommand{\ojoin}{\setbox0=\hbox{$\bowtie$}%
    \rule[-.02ex]{.25em}{.4pt}\llap{\rule[\ht0]{.25em}{.4pt}}}
\newcommand{\insertinnerjoin}{\mathbin{\bowtie\mkern-5.8mu}\;}
\newcommand{\insertleftouterjoin}{\mathbin{\ojoin\mkern-5.8mu\bowtie}}
\newcommand{\insertrightouterjoin}{\mathbin{\bowtie\mkern-5.8mu\ojoin}}
\newcommand{\insertfullouterjoin}{\mathbin{\ojoin\mkern-5.8mu\bowtie\mkern-5.8mu\ojoin}}


% Advanced Formats.
\newcommand{\insertimage}[3]
{
    \begin{figure}[H]
        \centering
        \includegraphics[width=#3\linewidth]{#1}
        \caption{#2}
        \label{fig:#2}
    \end{figure}
}

\newcommand{\insertalgorithm}[4]
{
    \begin{algorithm}[H]
        \caption{#1}
        \begin{algorithmic}
            \State \textbf{Input}: #2
            \State \textbf{Output}: #3
            \separate{5}#4
        \end{algorithmic}
        \label{alg:#1}
    \end{algorithm}
}

\newcommand{\inserttable}[4]
{
    \begin{table}[H]
        \centering
        \begin{tabular}{|#1|}
            \hline #3 \\ \hline #4 \\ \hline
        \end{tabular}
        \caption{#2}
        \label{tab:#2}
    \end{table}
}

%%%



%%% Document Setup

% Title Page.
\title{\thetitle\\[0.5em]\large{\thesubtitle}}
\author{\theauthor}
\date{\thedate} 


% Custom Page Header.
\fancyhead{}
\fancyhead[R]{\thetitle\text{ - }\thesubtitle}

% Custom Page Footer.
\fancyfoot{}
\fancyfoot[L]{Written by \theauthor \\ 
\textit{\LaTeX \,Template v.\templateversion: Copyright \date{\yeardate\today}, Matthew Emerson}} % Do not edit.
\fancyfoot[R]{Page \thepage \, of\, \pageref{LastPage}}

%%%



%%% Page Variables

\newcommand{\thetitle}{Mathematics for Computer Science} % Insert the document title here.
\newcommand{\thesubtitle}{COMP1021: Durham University} % Insert the document subtitle here.
\newcommand{\theauthor}{Matthew Emerson} % Insert the name(s) of the author(s) here.
\newcommand{\thedate}{\monthyeardate\today} % Replace this if you want a static date.

%%%



%%% Setting Up Pages.

\pagestyle{empty}\begin{document}
\maketitle\thispagestyle{empty} % Title Page.
\newpage\tableofcontents % Contents Page.

\clearpage\setcounter{page}{1}\newpage % Excludes the title and contents page from page count.
\pagestyle{fancy}

%%%



%%% Document Content

\part{Calculus}\label{part:1.calc}

\section{Introduction to Calculus}

\subsection{Sets}

\definition{1.1.1}{A \underline{set} is an unordered collection of distinct elements denoted as \(A = \{x | \textit{condition} \}\).}

\separate{5}Some common set operations include:
\begin{itemize}
    \item \definition{1.1.2}{The \underline{union} of two sets, \(A \cup B\), includes all the elements of both sets.}
    \item \definition{1.1.3}{The \underline{intersection} of two sets, \(A \cap B\), includes any elements which are present in both sets.}
    \item \definition{1.1.4}{If \(A\) and \(B\) are finite sets, then the \underline{principle of inclusion-exclusion} states that 
    \begin{equation}
        A \cup B = A + B - (A \cap B)
    \end{equation}}
    \item \definition{1.1.5}{The set \(A\) is a \underline{subset} of \(B\), \(A \subseteq B\), if every element of \(A\) is an element of \(B\). \(A\) is a \underline{proper subset} of \(B\), \(A \subset B\), if \(A\) is a subset of \(B\) and there is at least one element of \(B\) not in \(A\).}
    \item \definition{1.1.6}{The \underline{powerset} of a set \(A\) is the set of all subsets of \(A\). 
    \begin{equation}
        \mathcal{P}(A) = \{A'|A' \subseteq A\}
    \end{equation}}
    \item \definition{1.1.7}{The \underline{cartesian product} of two sets, \(A \times B\), is the set of all possible ordered pairs with the first element from \(A\) and the second from \(B\).}
\end{itemize}

\separate{5}\definition{1.1.8}{For a finite set, the \underline{cardinality} is the number of distinct elements that it contains.}

\subsection{Functions}

\separate{5}For a function \(f: A \rightarrow B\),
\begin{itemize}
    \item \definition{1.2.1}{\(A\) is known as the \underline{domain} of \(f\).}
    \item \definition{1.2.2}{\(B\) is known as the \underline{codomain} of \(f\).}
    \item \definition{1.2.3}{If \(f(a)=b\), then \(b\) is the \underline{image} of \(a\).}
    \item \definition{1.2.4}{The \underline{pre-image} of \(b\) is the set \(\{a|f(a)=b\}\).}
    \item \definition{1.2.5}{\(f\) forms an \underline{injection} if every element of \(A\) maps to a unique element in \(B\), even if some elements of \(B\) don't have a mapping.}
    \item \definition{1.2.6}{\(f\) forms a \underline{surjection} if every element of \(B\) maps to at least one element of \(A\).}
    \item \definition{1.2.7}{\(f\) forms a \underline{bijection} if it is both injective and surjective.}
\end{itemize}

\separate{5}\definition{1.2.8}{If for an unbounded set \(A\), there is a bijective function \(f: A \rightarrow \N\), then we say \(A\) is \underline{countably infinite}, otherwise \(A\) is known as \underline{uncountably infinite}.}

\subsection{Hyperbolic Functions}

\definition{1.3.1}{Given the unit hyperbola \(x^2-y^2=1\), the \underline{hyperbolic functions} are defined by the coordinates of a point at which area \(\frac{1}{2}\alpha\) has been swept out.}
\begin{equation*}
    sinh(\alpha) = \frac{e^\alpha - e^{-\alpha}}{2} \qquad cosh(\alpha) = \frac{e^\alpha + e^{-\alpha}}{2} \qquad tanh(\alpha) = \frac{e^x - e^{-x}}{e^x + e^{-x}} \qquad \sigma(\alpha) = \frac{1}{1 + e^{-\alpha}}
\end{equation*}
\begin{equation*}
    sech(\alpha) = \frac{2}{e^\alpha + e^{-\alpha}} \qquad cosech(\alpha) = \frac{2}{e^x - e^{-x}} \qquad coth(\alpha) = \frac{e^\alpha + e^{-\alpha}}{e^\alpha - e^{-\alpha}}
\end{equation*}
These functions can then be differentiated as follows:
\begin{equation*}
    \frac{d}{dx}\left[sinh(\alpha)\right] = cosh(\alpha) \qquad\qquad \frac{d}{dx}\left[cosh(\alpha)\right] = sinh(\alpha)
\end{equation*}
\begin{equation*}
    \frac{d}{dx}\left[tanh(\alpha)\right] = sech^2(\alpha) \qquad\qquad \frac{d}{dx}\left[\sigma(\alpha)\right] = \sigma(\alpha)(1-\sigma(\alpha))
\end{equation*}

\section{Multivariate Functions}

\subsection{Multivariate Derivatives}

\definition{2.1.1}{For a function \(f(x,y)\), if we take \(f=f_0\) to be fixed but unknown, then we can differentiate \(f\) as a function of \(x\) at \(x=x_0\), giving us the \underline{partial derivative of \(f\) with respect to \(x\)} which gives the gradient when moving in the \(x\)-direction from any point.}
\begin{equation}
    \frac{\delta f}{\delta x} \left[f(x,y)\right] = \lim_{x \rightarrow x_0} \left(\frac{f(x,y_0)-f(x_0,y_0)}{x-x_0}\right) = f_x
\end{equation}

\separate{5}\definition{2.1.2}{The \underline{vector of partial derivatives}, denoted as \(\nabla f\), contains all the partial derivatives of a function. This forms a vector pointing in the direction of the greatest rate of increase of \(f\), where its magnitude is the rate of increase in that direction. For example, for a function \(f(x,y)\), 
\begin{equation} 
    \nabla f = \begin{pmatrix}f_x\\f_y\end{pmatrix}
\end{equation}}

\separate{5}\definition{2.1.3}{The \underline{directional derivative} of the function \(f(x,y)\) gives the gradient at any point whilst looking in the direction given by the unit vector \(\mathbf{v} = \begin{pmatrix}v_x\\v_y\end{pmatrix}\).}
\begin{equation}
    \nabla_{\mathbf{v}} f = \lim_{h \rightarrow 0}\frac{f(x_0 + h\mathbf{v}) - f(x_0)}{h}
\end{equation}
If the slope is approximately flat at a small scale, then the slope can be determined by
\begin{equation}
    \nabla_{\mathbf{v}} f = \mathbf{v} \cdot \nabla f
\end{equation}

\separate{5}From this, we get that, for a function \(f(\mathbf{x_0})\) and a direction \(\mathbf{v}\),
\begin{equation}
    f(\mathbf{x_0} + h \mathbf{v}) \approx f(\mathbf{x_0}) + h \mathbf{v} \cdot \nabla f(\mathbf{x_0})
\end{equation}

\newpage\subsection{Multivariate Extrema}

\definition{2.2.1}{A function that is continuous is in \underline{class \(C^0\)}, and a function that is continuous and has a continuous derivative is in \underline{class \(C^1\)}. A function \(f\) such that \(f',\;f'',\;...,\;f^{(k)}\) exist and are continuous is in \underline{class \(C^k\)}.}

\separate{5}\definition{2.2.2}{A function \(f\) such that derivatives of all orders exist and are continuous is in class \(C^\infty\) is also \underline{smooth}.}

\separate{5}\definition{2.2.3}{All the second-order partial derivatives of a function \(f(x,y)\) can be gathered into a \underline{Hessian matrix}.}
\begin{equation}
    \mathbf{H}_f = \begin{bmatrix} f_{xx} & f_{xy} \\ f_{yx} & f_{yy} \end{bmatrix}
\end{equation}
This can be used to determine the form of bivariate extrema. Suppose \(f(x,y) \in C^2\) and \(f_x(x_0,y_0) = 0 = f_y(x_0,y_0)\) then,
\begin{itemize}
    \item \((x_0,y_0)\) is a local maximum if \(det(\mathbf{H}_f) > 0\) and \(f_{xx} < 0\) at \((x_0,y_0)\).
    \item \((x_0,y_0)\) is a local minimum if \(det(\mathbf{H}_f) > 0\) and \(f_{xx} > 0\) at \((x_0,y_0)\).
    \item \((x_0,y_0)\) is a saddle point if \(det(\mathbf{H}_f) < 0\).
    \item If \(det(\mathbf{H}_f) = 0\), then the test is inconclusive.
\end{itemize}

\separate{5}At a stationary point \(x_0\), the quadratic approximation of a function \(f(x)\) is
\begin{equation}
    f(x_0 + \mathbf{v}) \approx f(x_0) + \frac{1}{2}\mathbf{v}^T H_f(x_0)\mathbf{v}
\end{equation}
But as \(\mathbf{H}_f(x_0)\) is real and symmetric, it has orthogonal eigenvectors \(\mathbf{e_1},\;\mathbf{e_2},\;...,\;\mathbf{e_n}\) with eigenvalues \(\lambda_1,\;\lambda_2,\;...,\;\lambda_n\), so we can write
\begin{equation}
    \mathbf{v} = c_1\mathbf{e_1} + c_2\mathbf{e_2} + ... + c_n\mathbf{e_n}
\end{equation}
When substituted back into the original equation, this gives that
\begin{equation}
    f(x_0 + \mathbf{v}) \approx f(x_0) + \frac{1}{2}\left(\lambda_1(c_1)^2 + \lambda_2(c_2)^2 + ... + \lambda_n(c_n)^2\right)
\end{equation}
Given this, we can use the eigenvalues to determine the nature of the stationary point \(x_0\),
\begin{itemize}
    \item If all eigenvalues are positive, then the stationary point is a local minimum.
    \item If all eigenvalues are negaitve, then the stationary point is a local maximum.
    \item If there is a mixture of positive and negative eigenvalues, then the stationary point is a saddle point.
\end{itemize}

\subsection{Vector-Valued Functions}

\definition{2.3.1}{A \underline{vector-valued} function takes an \(n\)-dimensional vector as input, and outputs an \(m\)-dimensional vector, e.g., \(f:\R^n\rightarrow\R^m\).}

\separate{5}\definition{2.3.2}{The \underline{Jacobian matrix} of a vector-valued function \(f:\R^n\rightarrow\R^m\) can be formed by setting \(\mathbf{J}_{i,j} = \frac{\delta f_i}{\delta x_j}\) to form,}
\begin{equation}
    \mathbf{J}_f = \begin{pmatrix}
        \frac{\delta f_1}{\delta x_1} & \dots & \frac{\delta f_1}{\delta x_n} \\
        \vdots & \ddots & \vdots \\
        \frac{\delta f_m}{\delta x_1} & \dots & \frac{\delta f_m}{\delta x_n} \\
    \end{pmatrix}
\end{equation}
For a multivariate function \(f:\R^n\rightarrow\R^m\), the Jacobian matrix can be used to linearly approximate \(f\) at \(x\),
\begin{equation}
    f(x + h\mathbf{v}) \approx f(x) + \mathbf{J}_f(x)(h\mathbf{v})
\end{equation} 

\newpage\part{Linear Algebra}\label{part:2.la}

\newpage\part{References}

\section{Module Professors}

The below is accurate as of the 2025/26 academic year and will not be updated in subsequent years.
\begin{itemize}
    \item \hyperref[part:1.calc]{\textit{Calculus}} is taught by Doctor Junyan Hu and Doctor Sukanya Pandey.
        
    \separate{3}Doctor Junyan Hu:
    \begin{itemize}
        \item \textbf{Office Location}: Room 2060, MCS Building
        \item \textbf{Preferred Contact Methods}: junyan.hu@durham.ac.uk
        \item \textbf{Pronouns}: he, him
    \end{itemize}
        Doctor Sukanya Pandey:
    \begin{itemize}
        \item \textbf{Office Location}: Room 2011, MCS Building
        \item \textbf{Preferred Contact Methods}: sukanya.pandey@durham.ac.uk
        \item \textbf{Pronouns}: he, him
    \end{itemize}

    \item \hyperref[part:2.la]{\textit{Linear Algebra}} is taught jointly by Professor Steven Bradley and Doctor Karl Southern.
    
    \separate{3}Professor Steven Bradley:
    \begin{itemize}
        \item \textbf{Office Location}: Room 2027, MCS Building
        \item \textbf{Preferred Contact Methods}: s.p.bradley@durham.ac.uk
        \item \textbf{Pronouns}: he, him
    \end{itemize}
    Doctor Karl Southern
    \begin{itemize}
        \item \textbf{Office Location}: Room 2007, MCS Building
        \item \textbf{Preferred Contact Methods}: karl.southern@durham.ac.uk
        \item \textbf{Pronouns}: he, him
    \end{itemize}
\end{itemize}

\section{Cited Works}

\end{document}

%%%